% !PRINTED 337
\weberpage{337}

Die Kongruenz (15) ist nach (12) und (7) gleichbedeutend
mit den $n^2$-Kongruenzen

\[
y_{h,k}=\overset{r}{\Sigma}\,x_{r}\,c_{h,k}^{(r)}\equiv 0\pmod{p} % Source scan prints r above Sigma and no lower index.
\]
und aus (6) folgt alsdann:
\[
x_s\equiv 0\pmod{p},
\]
und dies besagt, da\ss{} die Substitutionen (17) f\"ur jeden Modul $p$
linear unabh\"angig sind.

\gesperrt{Nach Gauss (Disq. aritm. art. 235) hei\ss t eine Form $T$
aus den beiden Formen $U, V$ komponiert oder zusammengesetzt, wenn $T$ durch eine bilineare Substitution f\"ur
die Variable $t$, deren Gleichungen f\"ur jeden Modul linear
unabh\"angig sind, in das Produkt $UV$ \"ubergeht.}

Danach haben wir den Satz:

6. \gesperrt{Die Form, die zu dem Produkt zweier Ideale $\mathfrak{a}, \mathfrak{b}$
geh\"ort, ist aus den Formen der Ideale $\mathfrak{a}$ und $\mathfrak{b}$
komponiert}\footnote{Die S\"atze, die \gesperrt{Gauss} an der erw\"ahnten Stelle durch sehr weitl\"aufige
Rechnung f\"ur bin\"are quadratische Formen beweist, sind hier in gr\"o\ss ter
Allgemeinheit durch einfache Betrachtungen abgeleitet. Vgl. \gesperrt{Dirichlet-Dedekind}, Vorlesungen \"uber Zahlentheorie, \S 182. F\"ur bin\"are quadratische
Formen: \gesperrt{Dedekind}, Crelles Journal, Bd. 129; \gesperrt{H. Weber}, G\"ottinger Nachrichten 1907.}.

\begin{center}\rule{0.2\linewidth}{0.4pt}\end{center}
